\chapter{Pipeline, từng dòng một}
\label{ch:pipeline}

Mọi thứ trong cuốn sách này hội tụ về một file Groovy 70 dòng. Một
Jenkinsfile dạng khai báo sống trong chính repository mà nó build ---
pipeline dưới dạng mã, có phiên bản và được xem xét như mọi thứ khác.

\section{Phần mở đầu và environment}

\begin{lstlisting}
pipeline {
  agent any                                  (*@\co{1}@*)
  options {
    timestamps()
    disableConcurrentBuilds()                (*@\co{2}@*)
    buildDiscarder(logRotator(numToKeepStr: '10'))
  }
  triggers { pollSCM('H/5 * * * *') }        (*@\co{3}@*)
  environment {
    JAVA_HOME = "${env.JDK25_HOME}"          (*@\co{4}@*)
    PATH      = "${env.JDK25_HOME}/bin:${env.PATH}"
    TAG       = "${env.GIT_COMMIT.take(7)}"  (*@\co{5}@*)
    NS        = 'ts'
    REGISTRY  = 'localhost:5000'
    SERVICES  = 'config-server discovery-server course-service api-gateway auth-service dictionary-service notification-service'
    MODULES   = 'common,config-server,discovery-server,course-service,api-gateway,auth-service,dictionary-service,notification-service'
  }
\end{lstlisting}

\begin{description}
  \item[\co{1}] Bất kỳ executor nào trên controller --- thỏa hiệp
  homelab của Chương~22 gói trong một từ.
  \item[\co{2}] Hai build cùng lúc trên hai nhân sẽ giẫm đạp lên
  nhau; các lần chạy xếp hàng cũng được gộp lại (một commit mới trong
  lúc đang build thay thế những lần đang chờ).
  \item[\co{3}] Polling vì NAT chặn webhook. Chữ \code{H} là phân bố
  theo hash: ``một phút cố định nào đó trong mỗi 5 phút,'' để nhiều
  job không cùng poll vào phút :00.
  \item[\co{4}] Công tắc JDK theo từng pipeline mà image của
  controller đã chuẩn bị: build thấy Java\,25, bản thân Jenkins vẫn ở
  21.
  \item[\co{5}] Tag triển khai là SHA git rút gọn --- mọi image truy
  vết được về commit của nó, mọi lần deploy ghim vào một tag bất biến
  (bài học tag-so-với-digest của Chương~13 được áp dụng).
  \item[\code{SERVICES}/\code{MODULES}] sổ đăng ký của pipeline về
  những gì tồn tại: \code{MODULES} cấp cho Maven (phân cách bằng dấu
  phẩy, gồm cả \code{common} được build nhưng không sinh image);
  \code{SERVICES} cấp cho kubectl (phân cách bằng dấu cách, chỉ những
  thứ triển khai được). Thêm một service vào cluster nghĩa là mở rộng
  cả hai --- bước thủ công duy nhất cho mỗi service mới.
\end{description}

\section{Ba stage}

\begin{lstlisting}
stage('Build') {
  steps {
    dir('Project/TeacherSupporter') {          (*@\co{1}@*)
      sh 'mvn -B -DskipTests -pl ${MODULES} -am verify'  (*@\co{2}@*)
    } } }
stage('Push images') {
  steps {
    dir('Project/TeacherSupporter') {
      sh 'mvn -B -DskipTests -pl ${MODULES} -am package jib:build -Djib.to.tags=${TAG}'  (*@\co{3}@*)
    } } }
stage('Deploy') {
  steps {
    dir('Project/TeacherSupporter') {
      withCredentials([file(credentialsId: 'kubeconfig-ts',
                            variable: 'KUBECONFIG')]) {   (*@\co{4}@*)
        sh '''
          set -eu                                         (*@\co{5}@*)
          kubectl apply -k deploy/k8s/base                (*@\co{6}@*)
          for s in ${SERVICES}; do
            kubectl -n ${NS} set image deployment/${s} \
                ${s}=${REGISTRY}/${s}:${TAG}              (*@\co{7}@*)
          done
          for s in ${SERVICES}; do
            kubectl -n ${NS} rollout status \
                deployment/${s} --timeout=300s            (*@\co{8}@*)
          done
        '''
      } } } }
\end{lstlisting}

\begin{description}
  \item[\co{1}] Gốc repository là một monorepo học tập; dự án nằm
  trong một thư mục con, nên mọi stage đều \code{dir()} vào đó.
  \item[\co{2}] \code{-pl} (project list) + \code{-am} (also make
  dependencies): chỉ build các module được liệt kê và những gì chúng
  cần --- kinh tế học monorepo trên hai nhân. \code{-B} là chế độ
  batch (không có đầu ra tương tác). \code{-DskipTests} là sự làm yếu
  cổng chặn được chấp nhận có ý thức từ hộp cạm bẫy của Chương~21.
  \item[\co{3}] Một lệnh build và push cả bảy image --- Chương~13 gói
  trong một dòng. \code{-Djib.to.tags} thêm tag SHA bên cạnh tag
  SNAPSHOT. Chữ \code{package} đứng trước \code{jib:build} không phải
  để trang trí --- xem hộp cạm bẫy bên dưới; nó đã tốn hai lần build
  thất bại trong ngày cột mốc 5 và 6 ra lò.
  \item[\co{4}] Credential dạng file hiện hình thành một file tạm;
  \code{KUBECONFIG} là biến môi trường mà kubectl đọc. Ngoài khối này
  file không tồn tại; trong log nội dung của nó không bao giờ xuất
  hiện.
  \item[\co{5}] Thất bại khi gặp bất kỳ lỗi nào và khi gặp biến chưa
  đặt --- thiếu nó, một lần apply thất bại sẽ bị lướt qua và
  ``đã triển khai'' được báo cáo.
  \item[\co{6}] Hội tụ cluster về các manifest --- lũy đẳng
  (Chương~14), nên chạy lại an toàn ở mỗi build dù có gì thay đổi hay
  không.
  \item[\co{7}] Trỏ lại từng Deployment vào image tag theo SHA. Đây
  mới là thứ thật sự kích hoạt rollout --- và lưu ý bản thân manifest
  vẫn giữ tag SNAPSHOT, nên chỉ \emph{apply} thôi sẽ không triển khai
  lại; việc ghim của pipeline và giá trị mặc định của manifest cùng
  tồn tại.
  \item[\co{8}] Cổng kiểm chứng: chặn cho đến khi mỗi rollout hoàn
  tất; làm build thất bại nếu có pod nào không bao giờ Ready trong
  vòng năm phút. Dòng này là thứ khiến màu xanh của pipeline có nghĩa
  ``đang chạy,'' chứ không phải ``các lệnh đã thực thi'' --- nó ủy
  thác sự thật cho các probe của Chương~17.
\end{description}

\begin{pitfall}[Could not find artifact com.ts:common:jar:1.0.0-SNAPSHOT]
Triệu chứng: stage Build qua, rồi stage Push thất bại ở module đầu
tiên phụ thuộc vào \code{common}, không phân giải được một file jar mà
stage trước vừa build xong. Nguyên nhân: \code{mvn jib:build} là một
\emph{goal trần} --- Maven chạy thẳng nó, không có pha vòng đời nào,
nên không có gì trong reactor được biên dịch hay đóng gói trong phiên
đó, và các phụ thuộc anh em quay về kho cục bộ, nơi \code{common}
chưa bao giờ được cài (\code{verify} dừng trước \code{install}). Các
module không có phụ thuộc anh em (config-server, discovery-server)
qua trót lọt, khiến lỗi trông như đặc thù của từng module. Cách sửa:
đặt một pha lên trước --- \code{package jib:build} --- để reactor
đóng gói \code{common} trước khi Jib hỏi đến nó. Quy tắc chung: một
goal của plugin cần artifact của module khác phải đi cùng một pha
vòng đời, không được gọi một mình.
\end{pitfall}

Hai vòng lặp thay vì một là một điểm đúng đắn nhỏ: mọi rollout khởi
động trước, rồi mới đợi tất cả --- nếu không, việc chờ service N sẽ
tuần tự hóa việc khởi động service N+1, kéo dài tổng thời gian triển
khai mà chẳng được gì.

\section{Xử lý thất bại}

\begin{lstlisting}
post {
  success { echo "Deployed ${TAG} to namespace ${NS}." }
  failure {
    dir('Project/TeacherSupporter') {
      withCredentials([...]) {
        sh 'kubectl -n ${NS} get pods || true'
      } } } }
\end{lstlisting}

Khi thất bại, log build kết thúc bằng một danh sách pod --- lệnh đầu
tiên mà một con người sẽ chạy, đã được chạy sẵn. Phần \code{|| true}
giữ cho chính bước chẩn đoán không che mất trạng thái thoát thật của
build.

\section{Cùng pipeline đó, ở nơi khác}

Để hiệu chỉnh, đây là hình dạng y hệt dưới dạng GitHub Actions:

\begin{lstlisting}[language=yaml]
on: { push: { branches: [main] } }        # webhook, not polling
jobs:
  deploy:
    runs-on: ubuntu-latest                # disposable agent, not controller
    steps:
      - uses: actions/checkout@v4
      - uses: actions/setup-java@v4
        with: { distribution: temurin, java-version: 25 }
      - run: mvn -B -pl $MODULES -am verify
      - run: mvn -B -pl $MODULES -am jib:build -Djib.to.tags=${GITHUB_SHA::7}
      - run: kubectl apply -k deploy/k8s/base && ...
        env: { KUBECONFIG_DATA: ${{ secrets.KUBECONFIG }} }
\end{lstlisting}

Cùng các stage, cùng cách đánh tag, cùng ý tưởng gắn secret. Khác
biệt nằm ở sự đánh đổi hạ tầng: runner của GitHub là các VM dùng xong
bỏ (agent làm đúng cách, miễn phí) nhưng chúng sống trên internet ---
nên biến thể này sẽ cần registry và cluster truy cập được từ bên
ngoài, đúng ràng buộc đã định hình thiết kế homelab. Pipeline thì
khả chuyển; \emph{topology mạng mới là thứ thật sự khác nhau giữa các
nơi}.

\begin{quiz}
\begin{enumerate}
  \item Vì sao tag các lần deploy bằng SHA git khi tag SNAPSHOT cũng
  tồn tại? Nêu chế độ lỗi mà deploy chỉ-SNAPSHOT mời gọi (Chương~13)
  và khả năng truy vết mà SHA bổ sung.
  \item Giải thích \code{-pl ... -am} và nó tiết kiệm gì trong một
  monorepo. Điều gì hỏng nếu bạn quên mở rộng \code{MODULES} cho một
  service mới?
  \item Dòng đơn lẻ nào khiến thành công của pipeline này có nghĩa
  ``các pod đã Ready,'' và nó dựa vào bộ máy của chương nào?
  \item Vì sao dùng hai vòng lặp trong stage deploy thay vì
  set-image-rồi-đợi cho từng service?
  \item Liệt kê ba thứ thay đổi và ba thứ không đổi khi Jenkinsfile
  này được chuyển sang GitHub Actions.
\end{enumerate}
\end{quiz}
